% Options for packages loaded elsewhere
\PassOptionsToPackage{unicode}{hyperref}
\PassOptionsToPackage{hyphens}{url}
\PassOptionsToPackage{dvipsnames,svgnames,x11names}{xcolor}
%
\documentclass[
  11pt,
]{article}
\usepackage{amsmath,amssymb}
\usepackage{iftex}
\ifPDFTeX
  \usepackage[T1]{fontenc}
  \usepackage[utf8]{inputenc}
  \usepackage{textcomp} % provide euro and other symbols
\else % if luatex or xetex
  \usepackage{unicode-math} % this also loads fontspec
  \defaultfontfeatures{Scale=MatchLowercase}
  \defaultfontfeatures[\rmfamily]{Ligatures=TeX,Scale=1}
\fi
\usepackage{lmodern}
\ifPDFTeX\else
  % xetex/luatex font selection
\fi
% Use upquote if available, for straight quotes in verbatim environments
\IfFileExists{upquote.sty}{\usepackage{upquote}}{}
\IfFileExists{microtype.sty}{% use microtype if available
  \usepackage[]{microtype}
  \UseMicrotypeSet[protrusion]{basicmath} % disable protrusion for tt fonts
}{}
\makeatletter
\@ifundefined{KOMAClassName}{% if non-KOMA class
  \IfFileExists{parskip.sty}{%
    \usepackage{parskip}
  }{% else
    \setlength{\parindent}{0pt}
    \setlength{\parskip}{6pt plus 2pt minus 1pt}}
}{% if KOMA class
  \KOMAoptions{parskip=half}}
\makeatother
\usepackage{xcolor}
\usepackage[margin=1in]{geometry}
\setlength{\emergencystretch}{3em} % prevent overfull lines
\providecommand{\tightlist}{%
  \setlength{\itemsep}{0pt}\setlength{\parskip}{0pt}}
\setcounter{secnumdepth}{5}
\usepackage{microtype,amsmath,amssymb,booktabs}
\setlength{\parindent}{0pt}
\setlength{\parskip}{5pt plus 1pt minus 1pt}
\setcounter{secnumdepth}{2}
\providecommand{\tightlist}{\setlength{\itemsep}{0pt}\setlength{\parskip}{0pt}}
\ifLuaTeX
  \usepackage{selnolig}  % disable illegal ligatures
\fi
\IfFileExists{bookmark.sty}{\usepackage{bookmark}}{\usepackage{hyperref}}
\IfFileExists{xurl.sty}{\usepackage{xurl}}{} % add URL line breaks if available
\urlstyle{same}
\hypersetup{
  pdftitle={A transition with renters and owners},
  colorlinks=true,
  linkcolor={Maroon},
  filecolor={Maroon},
  citecolor={Blue},
  urlcolor={Blue},
  pdfcreator={LaTeX via pandoc}}

\title{A transition with renters and owners}
\usepackage{etoolbox}
\makeatletter
\providecommand{\subtitle}[1]{% add subtitle to \maketitle
  \apptocmd{\@title}{\par {\large #1 \par}}{}{}
}
\makeatother
\subtitle{Supporting proof for the illustrative model}
\author{}
\date{}

\begin{document}
\begin{center}
{\Large A transition with renters and owners}\par
\smallskip
{\small Supporting proof for the illustrative model}
\end{center}

\section{Result and scope}\label{result-and-scope}

There is an open set of economies with endogenous ownership, substantial
renting, positive child goods costs, and a positive rebated property tax
in which a small permanent increase in the financed share has a locally
unique converging competitive-equilibrium transition. Initial mean
fertility rises and the final population is larger. The original initial
old keep their pre-reform financial assets and housing titles.

The construction below has an owner share of \(11/21\), so almost half
the households rent. It proves local existence around an interior mixed
economy, not just continuity from a zero-renter limit. The neighborhood
and the permitted reform size are not given numerical lower bounds. It
does not prove global convergence, arbitrary changes of binding
constraints, or a general fertility sign. Its finite small reforms
eventually oscillate around the new steady state: cumulative fertility
can rise without fertility increasing at every date. The short
direct-allocation welfare result remains separate.

\section{Original choices and the maintained
branch}\label{original-choices-and-the-maintained-branch}

Keep the original flow utilities and value functions. Write
\(x=c-\chi n\), \(s=h-\kappa n\), and \(u_t=qr_t\) only to shorten the
proof. The primitives and the fertility weight \(\vartheta\) are
constant after the reform. There is no capital-gains tax. Define \[
 \rho_O=1+\beta(1+\gamma+\omega_B),\qquad
 \rho_R=1+\beta(1+\omega_B),\qquad a=h_R^{\max}.
\] We maintain the following uniformly strict conditional branches.
Young owners are at their down-payment limit, below their physical owner
cap, and save positively. Young and old renters are at their rental caps
and young renters save positively. Old owners are interior, with slack
retention and estate bounds. Adult goods, adult space, fertility and
estates are positive. Every conditional young owned size exceeds \(a\).
Conditional choices are evaluated in both tenures even when one is
selected less often.

Let \(F\) be the fixed entrant distribution of income and wealth, with
compact support and uniform margins for these inequalities. Additive
ownership tastes have a finite positive logistic scale \(\sigma_\xi\).
Conditional policies are independent of that taste draw; ownership
probabilities are not. The same probability must enter housing, births
and subsequent old-age demand.

At date \(t\), set \(w_i=y_i+b_i+T_t+qT_{t+1}\). The original dated
budgets and first-order conditions imply \[
 h_i^O=\frac{b_i}{(1-\phi)P_t},\quad h_i^R=a,\qquad
 \rho_Ox_i^O+\chi n_i^O=w_i-u_th_i^O,
\] \[
 \rho_Rx_i^R+\chi n_i^R=w_i-u_ta-qu_{t+1}a,\qquad
 \frac{\vartheta}{n_i^m}
 =\frac{\chi}{x_i^m}+\frac{\alpha\kappa}{h_i^m-\kappa n_i^m}.
\] Old choices for these young households satisfy \[
 c_{t+1,i}^{2,m}=\frac{\beta x_{t,i}^m}{q},\quad
 h_{t+1,i}^{2,O}=\frac{\beta\gamma x_{t,i}^O}{q u_{t+1}},\quad
 h_{t+1,i}^{2,R}=a,\quad
 e_{t+1,i}^m=\frac{\beta\omega_Bx_{t,i}^m}{q^2}.
\] Saving and old net assets are reconstructed from the original young
budget and mortgage repayment. In particular the estate is dated in
next-period goods; the denominator above is \(q^2\).

\section{Six variables suffice for aggregate
equilibrium}\label{six-variables-suffice-for-aggregate-equilibrium}

Use \(Z_t=(P_t,u_t,Y_t,O_t,M_t,R_t)\). The variables \(Y_t,O_t\) are
cohort masses. The quantity \(M_t/u_t\) is total old-owner housing;
\(M_t\) is not financial wealth. The variable \(R_t\) is total
old-renter housing. For \(t\geq1\), \[
 M_t=\frac{\beta\gamma}{q}Y_{t-1}
       \int\pi_{t-1}^O(i)x_{t-1,i}^O\,dF,
 \qquad
 R_t=aY_{t-1}\int(1-\pi_{t-1}^O(i))\,dF.
\] The user-cost identity and balanced tax rebate give \[
 P_{t+1}=\frac{(1+q\tau^p)P_t-u_t}{q},\qquad
 T_t=\frac{q\tau^pP_t\bar H}{Y_t+O_t}.
\] Given \(Z_t\), take \(v=(u_{t+1},Y_{t+1})\) as two unknowns and set
\(T_{t+1}=q\tau^pP_{t+1}\bar H/(Y_{t+1}+Y_t)\). Evaluate the preceding
original conditional choices and the logistic tenure probability. Solve
\[
 F_1(Z_t,v)=Y_t\int[\pi_i h_i^O+(1-\pi_i)a]dF
                    +M_t/u_t+R_t-\bar H=0,
\] \[
 F_2(Z_t,v)=Y_{t+1}-\nu Y_t
                 \int[\pi_i n_i^O+(1-\pi_i)n_i^R]dF=0.
\] The exact update is \[
 Z_{t+1}=G(Z_t,v)=
 \left(P_{t+1},u_{t+1},Y_{t+1},Y_t,
 \frac{\beta\gamma}{q}Y_t\int\pi_i x_i^OdF,
 aY_t\int(1-\pi_i)dF\right).
\] Where \(F_v\) is nonsingular, the implicit function theorem gives a
smooth local map \(Z_{t+1}=g(Z_t;\phi)\). This is a representation of
the original equilibrium equations, including the future rental cost in
tenure choice. It is not an assumed adjustment rule.

\subsection{A simple sufficient condition for the forward
solve}\label{a-simple-sufficient-condition-for-the-forward-solve}

For homogeneous income and wealth, suppose also that each old owner
occupies at least \(a\). These conditions imply nonsingularity of
\(F_v\) even with a positive property tax; neither a small-tax nor a
small-child-cost condition is needed for this particular step.

To see this, housing clearing fixes the owner probability at the given
state: \[
 \pi=\frac{\bar H-M/u-R-Ya}{Y(h^O-a)}\in(0,1).
\] Write \(z=u_{t+1}\) and let \(w\) vary with the future rebate. The
envelope derivatives of \(\Delta=W^O-W^R\) are \[
 \Delta_z=-\frac{\beta\gamma}{z}+\frac{qa}{x_R}<0,
 \qquad \Delta_w=\frac1{x_O}-\frac1{x_R}.
\] The strict inequality is exactly the old renter's cap condition.
Holding \(\pi\) fixed therefore implies \[
 \frac{dz}{dw}=\frac{1/x_O-1/x_R}
                    {\beta\gamma/z-qa/x_R}.
\] At fixed housing, fertility is increasing in resources when
\(\chi>0\): \(n_{m,w}>0\). The renter's dependence on \(z\) is
\(n_{R,z}=-qa n_{R,w}\); the owner's is zero at fixed \(w\). Hence along
the fixed-tenure-probability equation, \[
 \frac{d\bar n}{dw}=\pi n_{O,w}
 +(1-\pi)n_{R,w}
       \frac{\beta\gamma/z-qa/x_O}{\beta\gamma/z-qa/x_R}\geq0.
\] The numerator in the fraction is nonnegative precisely when
\(h^{2,O}\geq a\). Meanwhile \(w\) is nonincreasing in \(Y_{t+1}\)
because the future tax revenue is rebated to \(Y_{t+1}+Y_t\) households.
After solving housing for \(z\), the derivative of the birth-equation
residual with respect to \(Y_{t+1}\) is
\(1-\nu Y_t(d\bar n/dw)(dw/dY_{t+1})\geq1\). This proves local
nonsingularity. At \(\chi=0\) the resource derivatives of fertility
vanish and the same conclusion follows. This argument does not assert
existence at distant states or global continuation across a branch
change.

\section{The actual initial old}\label{the-actual-initial-old}

Keep the complete pre-reform old distribution \(G_0\) as exogenous data.
Let \(B_0\) be its owner mass, \(A_0\) its aggregate owner net financial
assets after the original mortgage repayment, and \(H_0\) its aggregate
inherited owner title. The latter is inherited purchased housing, not
housing retained after the old reoptimize. Let
\(L=\gamma/(1+\gamma+\omega_B)\).

The four initial restrictions are \[
 Y=Y_0,\quad O=O_0,\quad R=a(O_0-B_0),\qquad
 M=L[A_0+PH_0+T(P,Y_0,O_0)B_0].
\] The initial owner coefficient varies with the surprising price and
rebate; fixing it would wrongly suppress revaluation of existing title.
These four aggregates suffice for date-zero housing clearing. The full
\(G_0\) is still needed to reconstruct and check each old household's
consumption and estate. Uniform slack at the initial equilibrium
preserves those inequalities under a small reform.

\section{A local nonlinear transition
theorem}\label{a-local-nonlinear-transition-theorem}

Suppose the local map and boundary operator are at least \(C^2\),
\(Z^*\) is a stationary equilibrium on this branch, \(F_v\) is
nonsingular, and \(J=g_Z(Z^*;\phi)\) has four eigenvalues strictly
inside and two strictly outside the unit circle, counted with
multiplicity. Let \(B\) be the derivative of the four initial
restrictions with respect to \(Z_0\). Suppose \(B\) restricted to the
four-dimensional stable generalized eigenspace is invertible. Then
sufficiently small permanent changes of \(\phi\), starting from the
original stationary population and claims, have a locally unique
equilibrium sequence converging exponentially to the nearby stationary
state.

\textbf{Proof.} Since \(1\) is not an eigenvalue, \((I-J)\) is
invertible and the stationary state is a smooth function of the reform.
Center the sequence on this new state. Split the linear recurrence into
its stable and unstable generalized eigenspaces. Propagate the stable
part forward and obtain its initial value from \(B\); solve the unstable
part backward using its inverse and the requirement that the sequence
converge. Choose a weight \(\eta<1\) above all stable moduli and above
the inverse unstable moduli. The resulting operator has a bounded
inverse on sequences with norm \(\sup_t\eta^{-t}\|v_t\|\). The nonlinear
residual and boundary conditions are smooth on a uniform neighborhood,
so the sequence implicit function theorem applies. A zero stable
eigenvalue is harmless: only the unstable block is inverted. Uniform
household margins and differentiation under the fixed compact-support
integral reconstruct the full original equilibrium. This proves
uniqueness among nearby convergent sequences, not global uniqueness of
equilibrium. \(\square\)

\section{An exact mixed equilibrium}\label{an-exact-mixed-equilibrium}

All the following primitive values are rational except the ownership
taste location, which is specified by the exact value functions below:
\[
 q=\tfrac12,\quad\phi=\tfrac45,\quad b=\tfrac15,\quad
 \alpha=\beta=\omega_B=\tfrac25,\quad\gamma=\tfrac3{10},
\] \[
 \chi=\tfrac3{20},\quad\kappa=\tfrac12,\quad
 \vartheta=\tfrac{141}{400},\quad\nu=2,\quad
 a=\tfrac14,\quad h_O^{\max}=2,\quad\sigma_\xi=1,
\] \[
 \tau^p=\tfrac{467}{9250},\quad
 y=\tfrac{3521349281}{1677802000},\quad
 \bar H=\tfrac{68104}{68019}.
\] Set the logistic taste location to
\(\bar\xi=\sigma_\xi\log(\pi/(1-\pi))-W^O+W^R\) at the allocations
below, where \(\pi=11/21\). This is a primitive specification of a
finite taste distribution, not an equilibrium-dependent policy rule.
Hold it fixed during the credit reform. This construction is an
analytical example, not a calibration.

There is an exact steady state with \[
 P=Y=O=1,\quad u=\tfrac{9717}{18500},\quad
 T=\tfrac{3975571}{314587875},\quad\pi=\tfrac{11}{21},
\] \[
 (x_O,h_O,n_O)=(1,1,\tfrac34),\qquad
 (x_R,h_R,n_R)=(\tfrac{99}{74},\tfrac14,\tfrac9{40}).
\] Mean fertility is exactly \(1/2\). Old owner housing is
\(1480/3239>a\), old owner consumption is \(4/5\), and its estate is
\(16/25\). The strict original constraints are checked in the companion
receipt. In particular the young owner's marginal value of space is
\(.64>u\), whereas the interior old owner's equals \(u\). Thus this same
equilibrium also satisfies the short housing misallocation proposition.

\subsection{Exact differentiation}\label{exact-differentiation}

Differentiate the two conditional resource and fertility equations. For
\(A_n=\vartheta/n^2+\alpha\kappa^2/s^2\), \[
 \begin{pmatrix}\rho_m&\chi\\-\chi/x_m^2&A_n\end{pmatrix}
 \binom{dx_m}{dn_m}
 =\binom{dw-h_mdu-u\,dh_m-\mathbf1_{m=R}qa\,du_{t+1}}
              {\alpha\kappa\,dh_m/s_m^2}.
\] The owner derivative is \(dh_O=h_Od\phi/(1-\phi)-h_OdP/P\) and
\(dh_R=0\). Envelope derivatives give \[
 dW_O=\frac{dw-h_Odu}{x_O}
 +(\alpha/s_O-u/x_O)dh_O-\frac{\beta\gamma}{u_{t+1}}du_{t+1},
 \quad dW_R=\frac{dw-a\,du-qa\,du_{t+1}}{x_R},
\] and \(d\pi=\pi(1-\pi)(dW_O-dW_R)/\sigma_\xi\). All coefficients at
the specified equilibrium are rational. Therefore \[
 J=G_Z-G_vF_v^{-1}F_Z,\qquad
 Q=g_\phi=G_\phi-G_vF_v^{-1}F_\phi
\] are exact rational matrices. No numerical differentiation is used in
their certificate. Separate complex-step differentiation of the original
household helper checks the derivation.

\subsection{Roots, initial conditions and
signs}\label{roots-initial-conditions-and-signs}

The characteristic polynomial is \(z p_5(z)\) up to a positive constant,
with \[
\begin{split}
p_5(z)={}&11624093714366716003479040894500z^5\\
&+258366763596378495144999871594252z^4\\
&-1176628664469803650647365064033229z^3\\
&+408376917290861767134959387855371z^2\\
&-199454780254172995433079109028654z\\
&+694085946210757308455013341344.
\end{split}
\] An exact Sturm count gives three real roots. Rational isolating
intervals place them respectively near \(-26.1500318\), \(.003504813\),
and \(3.5930940\). Vieta's formulas put the remaining pair's product
near \(.18132116\) and sum near \(.32660224\). Rational inequalities
prove that the pair is nonreal with modulus in \((.42,.44)\). There are
exactly four stable roots including zero, and two unstable roots.

To check the boundary without rounded complex eigenvectors, let \(p_6\)
be the monic characteristic polynomial. A polynomial row \(\ell(z)\)
constructed from \(J\) satisfies exactly \[
 \ell(z)(J-zI)=-p_6(z)e_P^\top.
\] It is nonzero at each unstable root. Those two left eigenvectors
annihilate the stable generalized space. Rational interval Gaussian
elimination proves that the determinant of the matrix whose rows are
\(B\), \(\ell(r_-)\) and \(\ell(r_+)\) excludes zero. After the
normalization specified in the accompanying verification file, its value
is near \(.2580033\). This verifies the theorem's boundary condition.

The stationary derivative \(Z_\phi^*=(I-J)^{-1}Q\) has \[
 \frac{dY^*}{d\phi}
 =\frac{38021133978224611827644635285}
 {35866089542871200498536115024}>0.
\] For the initial derivative \(dZ_0\), solve \(BdZ_0=0\) and
\(\ell(r_\pm)dZ_0=\ell(r_\pm)Z_\phi^*\). Then \(dZ_1=JdZ_0+Q\). A
rational enclosure proves \[
 .0288<\frac{d\bar n_0}{d\phi}<.0289,\qquad
 1.0600<\frac{dY^*}{d\phi}<1.0602.
\] The full exact enclosures are stored with the accompanying
verification file receipt; the decimal bounds here are outward bounds
for readability. Smoothness gives both signs for sufficiently small
finite positive reforms. Strict margins, root counts, boundary
transversality and strict response signs also persist in an open
neighborhood of the primitives. Small nondegenerate compact-support
perturbations of entrant income and wealth are included, retaining the
uniform conditional branches. This is not a certified rectangular
parameter box or a claim about arbitrary heterogeneity.

\section{What happens after the initial fertility
rise}\label{what-happens-after-the-initial-fertility-rise}

The constructed transition need not be monotone. In fact every
sufficiently small positive reform eventually places fertility on both
sides of replacement, infinitely often; population similarly crosses its
new steady level infinitely often. The crossings decay exponentially.

Here is a finite-reform argument. At the baseline let
\(v=dZ_0-Z_\phi^*\), \(w_1=Jv\), and \(w_2=J^2v\). If the young-cohort
response had no component from the complex pair, then
\(e_Y^\top w_2-r_s e_Y^\top w_1=0\), where \(r_s\) is the small real
root; the zero mode has disappeared after one step. The rational
interval check instead places this expression near \(-.1346566\),
strictly below zero. Thus the complex modes occur in the
young-population response.

For a small reform \(\delta\), write the deviations from its new steady
state in smooth local stable coordinates. Let \(\lambda(\delta)\) be one
of the complex roots. Uniformly, \(|\lambda|\in(.42,.44)\) and the other
two stable roots have modulus below \(.01\).

\textbf{Uniform tail lemma.} The stable graph and its restricted map may
be chosen \(C^2\) jointly in the state and \(\delta\). At zero state its
nonlinear remainder and its state derivative are zero. On a sufficiently
small common neighborhood Taylor's theorem therefore bounds the
remainder by \(C\|v\|^2\), its state derivative by \(C\|v\|\), and its
parameter derivative by \(C\|v\|^2\). For the last bound one may use
\(C^3\) regularity; here the conditional formulas, implicit map and
parameter-dependent stable graph can be taken at least \(C^3\). The
fixed compact entrant support allows their first three derivatives to
pass under the integral. Choose a stable-coordinate norm whose linear
contraction is strictly below \(.5\). Shrinking the neighborhood
preserves a nonlinear contraction by \(\eta=.5\). The initial stable
coordinate is a \(C^2\) function of \(\delta\) and vanishes at zero, so
iteration gives \(\|v_t\|\le C|\delta|\eta^t\). Differentiating that
iteration, or using the weighted sequence IFT, gives
\(\|\partial_\delta v_t\|\le C\eta^t\). Consequently its nonlinear
forcing and derivative are bounded respectively by
\(C\delta^2\eta^{2t}\) and \(C|\delta|\eta^{2t}\). Constants are uniform
for all sufficiently small \(\delta\).

The parameter-dependent graph can also be obtained directly by the
one-sided sequence IFT, prescribing the stable initial coordinate and
leaving the unstable initial coordinate free. That construction requires
invertibility only on the unstable space, so the zero stable root causes
no problem here either. This supplies the regularity used in the lemma.
The complex coordinate therefore satisfies \[
 a_{t+1}=\lambda a_t+r_t,\qquad
 |r_t|\le C\delta^2\eta^{2t}.
\] Since \(\eta^2<|\lambda|\), the convergent sum \[
 A(\delta)=a_0+\sum_{s=0}^\infty\lambda^{-s-1}r_s
\] gives \(a_t=A(\delta)\lambda^t+O(\delta^2\eta^{2t})\). The fast
stable coordinates are \(O(|\delta|(.01)^t+\delta^2\eta^{2t})\);
unstable coordinates on the local stable graph are quadratic in the
stable coordinates. Consequently \[
 Y_t-Y^*(\delta)=2\operatorname{Re}
       [C_Y(\delta)\lambda(\delta)^t]
       +O(|\delta|(.01)^t+\delta^2\eta^{2t}).
\] The preceding nonzero linear signal implies \(C_Y'(0)\ne0\), so
\(C_Y(\delta)\ne0\) for every sufficiently small nonzero \(\delta\).
This differentiability follows from the lemma and differentiation of the
uniform convergent sum. Differentiating \(\lambda^{-s-1}\) adds a factor
proportional to \(s+1\), which remains summable because
\(\eta^2/|\lambda|<1\). The nonlinear summands sum to \(O(\delta^2)\).

A nonzero real projection of a rotating nonreal eigenmode has infinitely
many strictly positive and negative values bounded away from zero after
division by its modulus. For a rational rotation this follows over a
period; for an irrational rotation it follows from density on the
circle. The remainder above is smaller along both subsequences. Cohort
population therefore crosses its limiting level infinitely often. Total
household population has the same conclusion, since its leading
coefficient is multiplied by \(1+1/\lambda\ne0\). Finally, \[
 \bar n_t-\frac1\nu=
 \frac{Y_{t+1}-Y_t}{\nu Y_t},
\] whose leading complex coefficient is multiplied by
\((\lambda-1)/(\nu Y^*)\ne0\). Fertility too approaches replacement from
both sides. This result is compatible with higher initial fertility and
a larger final population. It makes no all-date welfare claim.

\section{A family of mixed economies}\label{a-family-of-mixed-economies}

The same stationary allocation can be supported at different taste
scales. Keep the listed primitives and stationary bundles, let
\(\sigma=\sigma_\xi>0\), and specify
\(\bar\xi(\sigma)=\sigma\log(11/10)-W^O+W^R\). This is a one-parameter
family of primitive taste distributions with the same stationary
allocation. Within each economy, both taste parameters are fixed during
the credit reform. Comparing different \(\sigma\) here is not a policy
experiment.

\textbf{Proposition.} Every member with \(1\leq\sigma\leq4\) has a
locally unique converging transition after a sufficiently small
permanent credit relaxation. Initial fertility and final population
rise. The result persists under small changes in the other primitives
and small entrant heterogeneity with uniform branch margins. The
constructed family's owner share is always \(11/21\); nearby economies
need not have that exact share. The oscillation theorem in Section 7 is
established near \(\sigma=1\), not on the entire interval. For every
\(\sigma>0\), the stationary young-cohort derivative is positive and
each tenure's stationary lifetime-value derivative is negative.

The stability part is algebraic for the entire positive half-line. Exact
differentiation gives \(J(\sigma)=J(0)+\sigma \mathbf a\mathbf b^\top\)
for two constant vectors \(\mathbf a,\mathbf b\), and \(Q(\sigma)\)
affine in \(\sigma\). Here \(J(0)\) is the algebraic limit of the
matrices; a degenerate taste distribution is not assumed in the economic
theorem. The nonzero-root polynomial is affine in \(\sigma\): \[
 p_{5,\sigma}(z)=p_{5,1}(z)+(\sigma-1)p_1(z),
\] where \(p_{5,1}\) is the integer polynomial in Section 6 and \[
\begin{split}
 p_1(z)={}&209842616525591903123960392500000z^4\\
 &-1024713381846718023989705716419600z^3\\
 &+426827306877372594433334439987600z^2\\
 &-162673088007827714095534626813600z\\
 &+495225810071736959616928665600.
\end{split}
\] Apply \(z=(1+w)/(1-w)\) to this quintic. The first column of the
Routh array for \((1-w)^5p_{5,\sigma}((1+w)/(1-w))\) has signs \[
 -\,,\quad-\,,\quad-\,,\quad+\,,\quad+\,,\quad-.
\] Each sign follows because the corresponding rational function of
\(\sigma\) has numerator coefficients all of the displayed sign and
denominator coefficients all positive. The exact coefficients are
provided with the reproducible proof. No row vanishes. Thus exactly two
roots have positive real part in \(w\), and three have negative real
part. There are two roots outside the unit circle in \(z\) and three
inside, in addition to the zero root of \(J\). The quintic's
discriminant is a degree-eight polynomial in \(\sigma\) with all
coefficients negative. Together with the three-real-root count at one
point, this proves three distinct real roots and one nonreal pair for
every \(\sigma>0\). The signs at \(z=-1,1\) and at infinity identify one
real root below \(-1\) and one above \(1\).

For the boundary and initial fertility sign, the rational interval check
covers the entire closed interval \([1,4]\) with 913 adjacent
subintervals. At each subinterval's endpoints, exact root isolation
brackets the two unstable real roots. A bracket whose polynomial signs
agree at both parameter endpoints brackets the root for every
intermediate parameter, since the polynomial is affine in \(\sigma\).
This is an interval proof, not a grid of unchecked intermediate
economies.

The rank-one representation provides a polynomial left eigenvector
independent of \(\sigma\). For \(A=J(0)\), take
\(\ell(z)=\mathbf b^\top\operatorname{adj}(zI-A)\). The exact identity
\[
 \ell(z)(J(\sigma)-zI)=-p_{6,\sigma}(z)\mathbf b^\top
\] follows from the determinant lemma and is also checked coefficient by
coefficient. Interval elimination then establishes a nonzero
initial-boundary determinant and a strictly positive initial fertility
derivative throughout each parameter interval. Adjacent endpoints and
full coverage are verified exactly. The family receipt gives uniform
positive lower bounds; the allowed finite reform size is still not
numerically quantified.

There is a simpler direct calculation for the stationary comparison. The
common denominator is positive: \[
D(\sigma)=29694400415045569770000\sigma+7922531612096952274579>0.
\] Exact differentiation yields \[
 \frac{dY^*}{d\phi}=
 \frac{55(625392361132072652215364400\sigma+
                 65900983926556653741810787)}{953456D(\sigma)}>0,
\] \[
 \frac{dW^{O*}}{d\phi}=
 -\frac{33(568192684411159099406576\sigma+
                 7091548916248680863487)}{224D(\sigma)}<0,
\] \[
 \frac{dW^{R*}}{d\phi}=
 -\frac{83313(49680061080395802000\sigma+
                 2808938751889938767)}{224D(\sigma)}<0.
\] These signs hold for every positive \(\sigma\). At \(\sigma=1\), the
value derivatives are approximately \(-2.25302\) for owners and
\(-.518979\) for renters. Since both conditional values fall and the
additive ownership taste is fixed, the maximum of the two values falls
for every fixed taste draw. This compares the same entrant in two steady
states. It does not rank different populations socially or describe all
transitional cohorts. In particular, the initial old may benefit from
the rise in the value of their inherited titles. The competitive credit
reform does not itself provide the compensation used in the
housing-allocation proposition.

The accompanying evidence index gives the exact calculations and
regeneration command. The checks include rational root and sign bounds,
complete interval coverage, and independent comparisons with the
original dated household equations. Finite numerical paths check those
equations; local convergence follows from the proof above. The earlier
all-owner transition and its primitive conditions remain in the longer
assessment.

\end{document}
